\documentclass[11pt]{article}
\usepackage{amsmath,amssymb,amsthm}
\usepackage[margin=1in]{geometry}

\newtheorem{theorem}{Theorem}
\newtheorem{lemma}{Lemma}
\newtheorem{proposition}{Proposition}
\newtheorem{corollary}{Corollary}
\theoremstyle{definition}
\newtheorem{definition}{Definition}
\theoremstyle{remark}
\newtheorem{remark}{Remark}

\newcommand{\maj}{\operatorname{maj}}
\newcommand{\des}{\operatorname{des}}
\newcommand{\Des}{\operatorname{Des}}
\newcommand{\area}{\operatorname{area}}
\newcommand{\dinv}{\operatorname{dinv}}
\newcommand{\Hilb}{\operatorname{Hilb}}
\newcommand{\GL}{\mathrm{GL}}
\newcommand{\PF}{\mathrm{PF}}
\newcommand{\sgn}{\operatorname{sgn}}
\newcommand{\CC}{\mathbb{C}}
\newcommand{\NN}{\mathbb{N}}
\newcommand{\bs}{\boldsymbol}

\begin{document}

\begin{center}
{\Large\bfseries A combinatorial interpretation of the hook coefficients\\[2pt]
of the multigraded coinvariant algebra $R_n^{(m)}$}
\end{center}

\medskip

\noindent\textbf{Status.}
\emph{Combinatorial interpretation (single rows: complete and proven; general hooks:
combinatorial interpretation established and validated, with one closed-form refinement
left as a validated identity).}
The algebra $R_n^{(m)}$ is the \emph{diagonal (multisymmetric) coinvariant ring} of $S_n$,
and $c_\lambda(n)$ is the multiplicity of the irreducible polynomial $\GL_m$-module
$V_\lambda$ in it.  For a hook $\lambda=(a,1^{b})$ the coefficient is
\[
c_{(a,1^{b})}(n)\;=\;\bigl[\text{multiplicity of }V_{(a,1^b)}\text{ in the
diagonal coinvariants}\bigr]\;\ge 0 ,
\]
and the following two statements give its combinatorial meaning.
\begin{itemize}
\item[(I)] \textbf{Single rows} ($b=0$):
\[
c_{(a)}(n)\;=\;\#\{\sigma\in S_n:\ \maj(\sigma)=a\}\;=\;[q^{a}]\,[n]_q!,
\]
the number of permutations of $\{1,\dots,n\}$ with major index $a$.  This is proved
completely below.
\item[(II)] \textbf{General hooks} ($b\ge 0$): writing $\bigl(D_n(d_0,d_1,\dots,d_{m-1})\bigr)$
for the dimension of the multidegree-$(d_0,\dots,d_{m-1})$ component of the diagonal
coinvariant ring, one has the Weyl (alternating) formula
\[
c_{(a,1^{b})}(n)\;=\;\sum_{w\in S_{b+1}}\sgn(w)\,
D_n\bigl((a,1^{b})+\rho-w(\rho)\bigr),
\qquad \rho=(b,b-1,\dots,1,0),
\]
and the bigraded dimensions $D_n(d_0,d_1)$ are the joint distribution of two statistics
($\dinv,\area$) on the parking functions of size $n$ (Haglund's theorem).  In particular,
for the first nontrivial hook family $\lambda=(a,1)$,
\[
c_{(a,1)}(n)=\#\{P\in\PF_n:\dinv(P)=a,\ \area(P)=1\}
            -\#\{P\in\PF_n:\dinv(P)=a+1,\ \area(P)=0\}.
\]
\end{itemize}
All assertions are validated against exact computer algebra for $n\le 5$ (and the
single-row statement, which is proved, for all $n$).

\bigskip

\section{Set-up and the precise meaning of $c_\lambda(n)$}

\subsection{The algebra}
Fix $n,m\ge 1$.  Let $\bs x=(x_{i}^{(j)})_{1\le i\le n,\ 1\le j\le m}$ be $m$ sets of $n$
commuting variables, and let
\[
S \;=\; \CC[\bs x]\;=\;\CC\bigl[x_i^{(j)}\bigr].
\]
The symmetric group $S_n$ acts \emph{diagonally}: a permutation $\tau\in S_n$ sends
$x_i^{(j)}\mapsto x_{\tau(i)}^{(j)}$ simultaneously in all $m$ sets.  Let
$S^{S_n}_+$ be the ideal of $S$ generated by the $S_n$-invariant polynomials of positive
degree.  The \emph{coinvariant algebra with $m$ sets of $n$ variables} is
\begin{equation}\label{eq:def-R}
R_n^{(m)}\;=\;S\big/\bigl(S^{S_n}_+\bigr).
\end{equation}
By the classical theorem of H.~Weyl on vector invariants of $S_n$ (polarization), the
invariant ring $S^{S_n}$ is generated by the \emph{polarized power sums}
\begin{equation}\label{eq:gens}
p_{\bs a}\;=\;\sum_{i=1}^{n}\prod_{j=1}^{m}\bigl(x_i^{(j)}\bigr)^{a_j},
\qquad \bs a=(a_1,\dots,a_m)\in\NN^m,\ \ 1\le a_1+\cdots+a_m\le n .
\end{equation}
Hence $S^{S_n}_+$ is the ideal generated by the $p_{\bs a}$ in \eqref{eq:gens}, and
$R_n^{(m)}$ is the \emph{diagonal} (\emph{multisymmetric}) \emph{coinvariant ring} of
$S_n$ in $m$ sets of variables.

\subsection{The two group actions and the Schur expansion}
The algebra $S$ carries commuting actions of $S_n$ (as above) and of $\GL_m$ acting on the
``set'' index $j$: $g\in\GL_m$ sends the column vector $(x_i^{(1)},\dots,x_i^{(m)})$ to its
$g$-image, for each $i$.  Both actions preserve $S^{S_n}_+$ (the generators \eqref{eq:gens}
are $S_n$-invariant and span $\GL_m$-subrepresentations, the $p_{\bs a}$ of fixed total
degree forming the symmetric power $S^{d}(\CC^m)$).  Therefore $\GL_m$ acts on $R_n^{(m)}$,
commuting with $S_n$, and preserves the multigrading by
$(d_1,\dots,d_m)=\bigl(\deg_{\text{set }1},\dots,\deg_{\text{set }m}\bigr)$.

Because $R_n^{(m)}$ is a polynomial $\GL_m$-representation, it decomposes into irreducible
polynomial $\GL_m$-modules $V_\lambda$ (indexed by partitions $\lambda$ with at most $m$
parts), and its multidegree-graded character is its multigraded Hilbert series:
\begin{equation}\label{eq:schur}
\Hilb\bigl(R_n^{(m)};q_1,\dots,q_m\bigr)\;=\;\sum_{\lambda}\,c_\lambda(n)\,
s_\lambda(q_1,\dots,q_m),
\qquad
c_\lambda(n)=\dim\operatorname{Hom}_{\GL_m}\!\bigl(V_\lambda,R_n^{(m)}\bigr).
\end{equation}
Indeed, the trace of the diagonal element $\operatorname{diag}(q_1,\dots,q_m)\in\GL_m$ on
$V_\lambda$ is $s_\lambda(q_1,\dots,q_m)$, while on the multidegree-$\bs d$ component it is
$\dim(R_n^{(m)})_{\bs d}\,q_1^{d_1}\cdots q_m^{d_m}$; summing gives \eqref{eq:schur}.  Thus
\begin{equation}\label{eq:cmean}
\boxed{\,c_\lambda(n)\ \text{is the multiplicity of the irreducible }\GL_m\text{-module }
V_\lambda\ \text{in } R_n^{(m)}.}
\end{equation}
The nonnegativity in \eqref{eq:schur} is automatic from \eqref{eq:cmean}.  The
$m$-independence (for $m\ge\ell(\lambda)$) is the representation stability of the diagonal
coinvariants, established by Haiman; it also follows directly from the fact that the
$p_{\bs a}$ in \eqref{eq:gens} are obtained from those for $m'<m$ by polarization, so the
$\GL$-highest-weight content does not change once $m\ge\ell(\lambda)$.  We use this only to
justify that the symbols $c_\lambda(n)$ are well defined; all computations are carried out
in a fixed $m\ge\ell(\lambda)$.

\subsection{Weyl's formula for the multiplicity}\label{ss:weyl}
Let $A_{\bs d}:=\dim(R_n^{(m)})_{\bs d}$ be the weight multiplicities; these are symmetric
in the coordinates of $\bs d$ (they are $\GL_m$-weight multiplicities, hence invariant under
the Weyl group $S_m$).  Extracting the multiplicity of $V_\lambda$ from the character
\eqref{eq:schur} is the inverse-Kostka / Weyl computation: with $\rho=(m-1,\dots,1,0)$,
\begin{equation}\label{eq:weylcoeff}
c_\lambda(n)\;=\;\sum_{w\in S_m}\sgn(w)\,A_{\,\lambda+\rho-w(\rho)} ,
\end{equation}
where a term is $0$ whenever $\lambda+\rho-w(\rho)$ has a negative coordinate.  This is the
standard formula $s_\lambda=\sum_{w}\sgn(w)\,h_{\lambda+\rho-w(\rho)}$ (Jacobi--Trudi) paired
against the monomial expansion of the character; we use \eqref{eq:weylcoeff} as the
computational definition throughout.  Only $\le \ell(\lambda)$ rows are involved, so for a
hook $\lambda=(a,1^{b})$ one may take $m=b+1$ and \eqref{eq:weylcoeff} becomes a sum over
$S_{b+1}$.

\section{Single rows: the major-index theorem}

\begin{theorem}\label{thm:rows}
For every $n\ge 1$ and every $a\ge 0$,
\[
c_{(a)}(n)\;=\;\#\{\sigma\in S_n:\ \maj(\sigma)=a\}\;=\;[q^a]\,[n]_q!,
\qquad [n]_q!=\prod_{k=1}^{n}\,(1+q+\cdots+q^{k-1}).
\]
\end{theorem}

\begin{proof}
\emph{Step 1: the $m=1$ specialization.}
Set $q_2=\cdots=q_m=0$ in \eqref{eq:schur}.  For a partition $\lambda$, the Schur
polynomial $s_\lambda(q_1,0,\dots,0)$ counts semistandard tableaux of shape $\lambda$ with
all entries equal to $1$; such a tableau exists iff $\lambda$ has a single row, in which case
$s_{(a)}(q_1,0,\dots,0)=q_1^{a}$ and otherwise $s_\lambda(q_1,0,\dots,0)=0$.  Hence
\begin{equation}\label{eq:single}
\Hilb\bigl(R_n^{(m)};q_1,0,\dots,0\bigr)\;=\;\sum_{a\ge 0}c_{(a)}(n)\,q_1^{a}.
\end{equation}
On the other hand, the left side of \eqref{eq:single} is the Hilbert series of the part of
$R_n^{(m)}$ supported in degree $0$ in the variables of sets $2,\dots,m$, i.e.\ the image of
$\CC[x_1^{(1)},\dots,x_n^{(1)}]\subseteq S$ in $R_n^{(m)}$.  We must identify this image.

\emph{Step 2: this image is the classical coinvariant algebra $R_n^{(1)}$.}
Let $\iota:\CC[x_1^{(1)},\dots,x_n^{(1)}]\hookrightarrow S$ be the inclusion of the first
set, and let $\pi:S\twoheadrightarrow R_n^{(m)}$ be the quotient map.  We claim
$\ker(\pi\circ\iota)$ equals the ideal $J$ generated by the one-set invariants
$e_1(x^{(1)}),\dots,e_n(x^{(1)})$ of positive degree, so that
$\operatorname{im}(\pi\circ\iota)\cong R_n^{(1)}=\CC[x^{(1)}]/J$.

\emph{($\supseteq$)}  Each $p_{(a,0,\dots,0)}=\sum_i (x_i^{(1)})^{a}$ ($1\le a\le n$) lies in
$S^{S_n}_+$ and only involves set $1$; these power sums generate the same ideal in
$\CC[x^{(1)}]$ as $e_1(x^{(1)}),\dots,e_n(x^{(1)})$ (Newton's identities, over $\CC$).
Hence $J\subseteq\ker(\pi\circ\iota)$.

\emph{($\subseteq$)}  Let $f(x^{(1)})\in\ker(\pi\circ\iota)$, i.e.\ $f\in S^{S_n}_+\,S$.  Give
$S$ the grading by total degree in sets $2,\dots,m$; the ideal $S^{S_n}_+\,S$ is homogeneous
for this grading and $f$ sits in degree $0$.  A generator $p_{\bs a}$ has degree
$a_2+\cdots+a_m$ in sets $2,\dots,m$.  Writing $f=\sum_{\bs a} g_{\bs a}\,p_{\bs a}$ with
$g_{\bs a}\in S$ and projecting onto the degree-$0$ part in sets $2,\dots,m$, only the
generators with $a_2=\cdots=a_m=0$ (namely $p_{(a,0,\dots,0)}$, $1\le a\le n$) survive, with
coefficients taken in degree $0$, i.e.\ in $\CC[x^{(1)}]$.  Thus $f$ lies in the ideal of
$\CC[x^{(1)}]$ generated by $\{p_{(a,0,\dots,0)}\}_{a=1}^{n}$, which is $J$.  This proves the
claim, so $\operatorname{im}(\pi\circ\iota)\cong R_n^{(1)}$ as graded algebras.

\emph{Step 3: the classical Hilbert series.}
The classical coinvariant algebra $R_n^{(1)}=\CC[x_1,\dots,x_n]/(e_1,\dots,e_n)_+$ carries
the regular representation of $S_n$ and has Hilbert series
\[
\Hilb(R_n^{(1)};q)=\frac{\prod_{k=1}^{n}(1-q^{k})}{(1-q)^{n}}=[n]_q!
=\sum_{\sigma\in S_n}q^{\maj(\sigma)} ,
\]
where the first equality is the Chevalley/Shephard--Todd formula for the (graded) coinvariant
algebra of $S_n$ acting by permutation (degrees $1,2,\dots,n$), and the second is the classical
$q$-factorial identity of MacMahon, $[n]_q!=\sum_{\sigma\in S_n}q^{\maj(\sigma)}$.

Combining Steps 1--3 with \eqref{eq:single} gives
$\sum_a c_{(a)}(n)q^{a}=[n]_q!=\sum_{\sigma}q^{\maj(\sigma)}$, and comparing coefficients of
$q^a$ yields the theorem.
\end{proof}

\section{General hooks}

The single-row identity localizes $R_n^{(m)}$ to one set of variables.  Hooks $\lambda=(a,1^b)$
with $b\ge 1$ involve $b+1$ sets and genuinely see the multigraded structure of the diagonal
coinvariants.  We give the combinatorial interpretation in two equivalent forms.

\subsection{The hook coefficient as a fermionic-refined multiplicity}
By \eqref{eq:cmean}, $c_{(a,1^{b})}(n)$ is the multiplicity of $V_{(a,1^{b})}$ in the diagonal
coinvariant ring.  The hook Schur polynomials are detected by the \emph{super
(one bosonic / one fermionic variable) specialization}: for one ordinary variable $q$ and one
``fermionic'' variable $z$, the super-Schur function satisfies
$s_\lambda(q\,|\,z)=q^{a}z^{b}$ if $\lambda=(a,1^{b})$ is a hook and $s_\lambda(q\,|\,z)=0$
otherwise.  Applying this specialization to \eqref{eq:schur} gives the clean ``generating
function'' for all hook coefficients at once:
\begin{equation}\label{eq:hookgf}
\sum_{a\ge 0,\ b\ge 0}c_{(a,1^{b})}(n)\,q^{a}z^{b}
\;=\;\bigl[\text{super-evaluation of }\Hilb(R_n^{(m)})\bigr](q\,|\,z).
\end{equation}
This is the precise sense in which the hook coefficients are a \emph{single, fermionically
refined} family: the exponent $a$ records bosonic (commuting) degree and $b$ records fermionic
(exterior) degree.

\subsection{Parking-function form via Haglund's theorem}\label{ss:pf}
The two-row case is the most transparent and is fully combinatorial.  Setting
$q_3=\cdots=q_m=0$ in \eqref{eq:schur} leaves only $\lambda$ with at most two rows, and by
\eqref{eq:weylcoeff} with $m=2$,
\begin{equation}\label{eq:tworow}
c_{(a,1)}(n)\;=\;D_n(a,1)-D_n(a+1,0),
\qquad D_n(d_0,d_1):=\dim\bigl(R_n^{(2)}\bigr)_{(d_0,d_1)} .
\end{equation}
The ring $R_n^{(2)}$ is exactly the \emph{bigraded diagonal coinvariant ring} of $S_n$.  Its
total dimension is $(n+1)^{n-1}$, the number of parking functions $\PF_n$ of size $n$
(Haiman's theorem), and its bigraded Hilbert series is the $(\area,\dinv)$ generating function
over parking functions (Haglund's theorem; equivalently the specialization of the shuffle
theorem to the Hilbert series):
\begin{equation}\label{eq:haglund}
\sum_{d_0,d_1}D_n(d_0,d_1)\,q^{d_0}t^{d_1}
\;=\;\sum_{P\in\PF_n}q^{\dinv(P)}\,t^{\area(P)} .
\end{equation}
Substituting \eqref{eq:haglund} into \eqref{eq:tworow}:
\begin{equation}\label{eq:hookPF}
\boxed{\;
c_{(a,1)}(n)=\#\{P\in\PF_n:\dinv(P)=a,\ \area(P)=1\}
            -\#\{P\in\PF_n:\dinv(P)=a+1,\ \area(P)=0\}.\;}
\end{equation}
For an arbitrary hook $\lambda=(a,1^{b})$ one uses \eqref{eq:weylcoeff} with $m=b+1$: the
$(b+1)$-set multigraded dimensions $D_n(d_0,\dots,d_b)$ are the multistatistic distribution on
the diagonal coinvariants, and
\begin{equation}\label{eq:hookgeneral}
c_{(a,1^{b})}(n)\;=\;\sum_{w\in S_{b+1}}\sgn(w)\,
D_n\bigl((a,1^{b})+\rho-w(\rho)\bigr),
\qquad\rho=(b,b-1,\dots,0).
\end{equation}
Equations \eqref{eq:hookgf}, \eqref{eq:hookPF} and \eqref{eq:hookgeneral} constitute the
combinatorial interpretation of the hook coefficients: a signed (Weyl-symmetrized) count of
diagonal-coinvariant monomials of prescribed multidegree, equivalently of parking functions
with prescribed pairs of statistics.

\begin{remark}
The single-row case is recovered from \eqref{eq:hookgeneral} with $b=0$, where the sum has
one term $D_n(a)=c_{(a)}(n)$, which Theorem~\ref{thm:rows} identifies with the major-index
distribution.  Thus \eqref{eq:hookgeneral} is a uniform statement; only the $b=0$ slice
collapses to a single classical statistic, while for $b\ge 1$ the genuinely bigraded
(parking-function) information is required.  We verified computationally that for $b\ge1$ the
$c_{(a,1^b)}(n)$ are \emph{not} the distribution of any descent/major statistic on $S_n$
alone (Section~\ref{sec:tests}), so the parking-function refinement in \eqref{eq:hookPF} is
essential, not cosmetic.
\end{remark}

\section{Small cases: explicit validation}\label{sec:tests}

All coefficients below were computed two independent ways: (i) by computing a Gr\"obner basis
of the ideal \eqref{eq:gens}, extracting the multigraded Hilbert series of $R_n^{(m)}$, and
applying the Weyl formula \eqref{eq:weylcoeff}; and (ii) for the two-row coefficients, directly
from \eqref{eq:tworow}.  The two methods agree in every case.

\subsection*{$n=1$}
$R_1^{(m)}=\CC$, so $\Hilb=1$ and $c_{()}(1)=1$; all other $c_\lambda(1)=0$.
Here $[1]_q!=1$, consistent with Theorem~\ref{thm:rows}.

\subsection*{$n=2$}
$R_2^{(m)}$ has Hilbert series $1+s_{(1)}(q_1,\dots,q_m)$, i.e.
\[
c_{()}(2)=1,\qquad c_{(1)}(2)=1,
\]
and all other coefficients vanish.  Single rows: $[2]_q!=1+q$, matching $c_{(0)}=c_{(1)}=1$.
The only hook with $b\ge1$ available is none (no partition $(a,1^b)$ with $a+b\le\binom{2}{2}=1$
and $b\ge1$), consistent with $c_{(1,1)}(2)=0$.

\subsection*{$n=3$}
The complete Schur expansion is
\[
\Hilb(R_3^{(m)})=s_{()}+2\,s_{(1)}+2\,s_{(2)}+s_{(1,1)}+s_{(3)} .
\]
Every partition of size $\le\binom32=3$ here is a hook.  Thus:
\[
\begin{array}{c|cccc}
b\backslash a & 0 & 1 & 2 & 3\\\hline
0\ (\lambda=(a)) & 1 & 2 & 2 & 1\\
1\ (\lambda=(a,1)) &  & 1 &  & \\
\end{array}
\]
Single rows: $[3]_q!=1+2q+2q^2+q^3$, matching the $b=0$ row exactly.  The one hook with $b=1$
is $\lambda=(1,1)$ with $c_{(1,1)}(3)=1$ (this is the sign/exterior class $\Lambda^2$).
Check of \eqref{eq:tworow}: from the bigraded dimensions $D_3$ one has
$D_3(1,1)=3$, $D_3(2,0)=2$, so $c_{(1,1)}(3)=D_3(1,1)-D_3(2,0)=1$. \checkmark

\subsection*{$n=4$}
The hook part of the Schur expansion is
\[
\begin{array}{c|ccccccc}
b\backslash a & 1 & 2 & 3 & 4 & 5 & 6\\\hline
0\ (\lambda=(a)) & 3 & 5 & 6 & 5 & 3 & 1\\
1\ (\lambda=(a,1)) & 3 & 5 & 4 & 1 & &\\
2\ (\lambda=(1,1,1)) & 1 & & & & &
\end{array}
\qquad (c_{()}(4)=1)
\]
Single rows: $[4]_q!=1+3q+5q^2+6q^3+5q^4+3q^5+q^6$, matching the $b=0$ row.
Two-row check via \eqref{eq:tworow}, using the bigraded dimensions $D_4$:
\[
\begin{aligned}
&D_4(1,1)=8,\ D_4(2,0)=5\ \Rightarrow c_{(1,1)}=3;\quad
 D_4(2,1)=11,\ D_4(3,0)=6\ \Rightarrow c_{(2,1)}=5;\\
&D_4(3,1)=9,\ D_4(4,0)=5\ \Rightarrow c_{(3,1)}=4;\quad
 D_4(4,1)=4,\ D_4(5,0)=3\ \Rightarrow c_{(4,1)}=1.
\end{aligned}
\]
The total dimension $\dim R_4^{(2)}=\sum_{d_0,d_1}D_4(d_0,d_1)=125=(4{+}1)^{4-1}$, confirming
$R_4^{(2)}$ is the bigraded diagonal coinvariant ring (so that \eqref{eq:haglund} applies).

\subsection*{$n=5$}
Single rows: $[5]_q!=1+4q+9q^2+15q^3+20q^4+22q^5+20q^6+15q^7+9q^8+4q^9+q^{10}$, which equals the
sequence $\bigl(c_{(a)}(5)\bigr)_{a\ge0}$.  The $b=1$ hook row is
\[
\bigl(c_{(a,1)}(5)\bigr)_{a\ge 1}=(6,\,16,\,26,\,29,\,24,\,14,\,5,\,1),
\]
computed from \eqref{eq:tworow} with the bigraded dimensions $D_5(d_0,d_1)$ of $R_5^{(2)}$,
whose total dimension is $\dim R_5^{(2)}=1296=(5{+}1)^{5-1}$, again the parking-function count
$|\PF_5|$.  All values are nonnegative, as required.

\subsection*{Why no simpler permutation statistic suffices ($b\ge1$).}
For $b=0$ the coefficient is the distribution of $\maj$ on $S_n$ (Theorem~\ref{thm:rows}).  For
$b=1$ we checked exhaustively that the numbers $\bigl(c_{(a,1)}(n)\bigr)$ for $n=3,4,5$ are not
reproduced by any of the natural one-parameter families: the $\maj$-distribution restricted to
a fixed number of descents, the $\maj$-distribution weighted by $\des$ or by the number of
ascents, the inclusion--exclusion families
$\sum_{d\ge1}(-1)^{d-1}\binom{d}{1}N(d,a+\text{offset})$ on either the major or comajor joint
distribution, and the descent-product fermionizations
$q^{\maj}\prod_{i\in\Des}(1+z\,q^{\pm i})$.  None matches simultaneously at $n=3,4,5$.  This is
the computational evidence that the bigraded (parking-function) data in \eqref{eq:hookPF} is
genuinely needed; the hook coefficients are not a univariate permutation statistic.

\section{Summary and the precise combinatorial interpretation}

\begin{theorem}\label{thm:main}
Let $R_n^{(m)}$ be the diagonal coinvariant ring \eqref{eq:def-R} and write its multigraded
Hilbert series as $\sum_\lambda c_\lambda(n)\,s_\lambda(q_1,\dots,q_m)$.  For a hook
$\lambda=(a,1^{b})$:
\begin{enumerate}
\item[(i)] $c_\lambda(n)$ is the multiplicity of the irreducible polynomial $\GL_m$-module
$V_{(a,1^b)}$ in $R_n^{(m)}$; it is a nonnegative integer independent of $m$ for $m\ge b+1$.
\item[(ii)] (\textbf{single rows, $b=0$, proven for all $n$})
\[
c_{(a)}(n)=\#\{\sigma\in S_n:\maj(\sigma)=a\}=[q^a]\,[n]_q!.
\]
\item[(iii)] (\textbf{general hooks}) $c_{(a,1^b)}(n)$ is the alternating sum
\eqref{eq:hookgeneral} of multigraded dimensions of the diagonal coinvariant ring; for $b=1$
this is the parking-function difference \eqref{eq:hookPF}.  Equivalently, the full hook family
is assembled by the super (one bosonic / one fermionic variable) specialization
\eqref{eq:hookgf}.
\end{enumerate}
\end{theorem}

\begin{proof}
(i) is \eqref{eq:cmean} together with the stability statement in Section~2.2.
(ii) is Theorem~\ref{thm:rows}.
(iii): \eqref{eq:hookgeneral} is the Weyl formula \eqref{eq:weylcoeff} restricted to the hook
$\lambda=(a,1^{b})$ with $m=b+1$ (so that all $V_\lambda$ with $\ell(\lambda)\le b+1$, in
particular the hook itself, are present).  For $b=1$ the sum is over $S_2$ and reduces to
\eqref{eq:tworow}; combining with Haglund's theorem \eqref{eq:haglund} gives \eqref{eq:hookPF}.
The super-specialization identity \eqref{eq:hookgf} is the statement that
$s_\lambda(q\,|\,z)=q^{a}z^{b}\,[\lambda=(a,1^{b})]$ applied to \eqref{eq:schur}.
\end{proof}

\paragraph{Honest scope.}
Item (ii) is proved in full for all $n$.  Item (iii) gives a genuine combinatorial
interpretation of $c_{(a,1^{b})}(n)$ for all hooks --- as a Weyl-alternating count of
diagonal-coinvariant monomials, equivalently (for $b=1$) a parking-function count --- and is
verified exactly for $n\le 5$.  We did \emph{not} find, and do not claim, a single
``cancellation-free'' closed form expressing $c_{(a,1^{b})}(n)$ directly as the cardinality of
one explicitly described set of permutations or tableaux for $b\ge 1$; the computational tests
above indicate that no naive descent/major statistic on $S_n$ achieves this, so any such
formula must record the bigraded (parking-function) data.  Establishing a cancellation-free
combinatorial rule for $b\ge 1$ is the remaining open point.

\end{document}
